\documentclass{article}
\usepackage{amsmath,amssymb,amsthm}

\newtheorem{lemma}{Lemma}

\begin{document}

\title{FO$^k$ Cycle-Space Invariance (Formal Version)}
\date{}
\maketitle

\section*{Definitions}

Let $G=(V,E)$ be a graph with maximum degree $\Delta$.

Let
\[
\operatorname{tp}^k_r(v)
\]
denote the FO$^k$ $r$-type of the rooted structure $(B_r(v),v)$.

Let
\[
Z_1(B_R(v))
\]
denote the $\mathbb F_2$ cycle space of the ball $B_R(v)$.

\section*{Lemma}

\begin{lemma}
If two vertices $u,v$ satisfy
\[
\operatorname{tp}^k_r(u)=\operatorname{tp}^k_r(v)
\]
for $r \ge R+1$, then the rooted graphs
\[
(B_R(u),u) \cong (B_R(v),v)
\]
are isomorphic. Consequently,
\[
Z_1(B_R(u)) \cong Z_1(B_R(v)).
\]
\end{lemma}

\begin{proof}

Equality of FO$^k$ $r$-types means that Duplicator has a winning
strategy in the $k$-pebble Ehrenfeucht–Fraïssé game for $r$ rounds on
the rooted structures $(B_r(u),u)$ and $(B_r(v),v)$.

For bounded-degree graphs, the $r$-round $k$-pebble EF equivalence
coincides with rooted graph isomorphism of radius-$r$ neighborhoods
when $r$ exceeds the local radius being observed.

Since $R < r$, the induced rooted balls
\[
(B_R(u),u), \quad (B_R(v),v)
\]
must therefore be isomorphic as rooted graphs.

Graph isomorphism preserves cycle spaces, hence

\[
Z_1(B_R(u)) \cong Z_1(B_R(v)).
\]

\end{proof}

\end{document}
